\documentclass{article}
\usepackage{graphicx}
\usepackage{xcolor}
\usepackage{float}
\usepackage{placeins}
\usepackage{rotating}
\usepackage{pdflscape}
\usepackage{caption}
\usepackage{subcaption}

\begin{document}
\section*{Figures and layout}

\begin{figure}[H]
\centering
\begin{subfigure}{0.45\textwidth}
  \centering
  \fcolorbox{blue!60!black}{blue!5}{\rule{0pt}{25mm}\rule{35mm}{0pt}}
  \caption{Synthetic panel A}
\end{subfigure}
\hfill
\begin{subfigure}{0.45\textwidth}
  \centering
  \rotatebox{8}{\fcolorbox{red!60!black}{red!5}{\rule{0pt}{25mm}\rule{35mm}{0pt}}}
  \caption{Synthetic panel B}
\end{subfigure}
\caption{A two-panel figure without external assets}
\end{figure}
\FloatBarrier

\begin{sidewaystable}
\centering
\caption{Rotated table}
\begin{tabular}{lrr}
Item & First & Second \\
A & 1 & 2 \\
B & 3 & 4 \\
\end{tabular}
\end{sidewaystable}

\begin{landscape}
This page exercises landscape PDF rotation metadata.
\end{landscape}
\end{document}
